\sectionguard
\subsection*{Charisms ordered to confession, service, and the common good\enspace\properrefs{Ep.}}

The appointed reading begins after Paul's heading concerning
\emph{spiritual things} (\latin{tōn pneumatikōn}, 1~Cor.~12:1), a phrase
whose range should not be reduced in advance to the later technical category
``charisms.'' Verses 2--3 first establish a passage from former pagan leading
toward mute idols (\latin{apagomenoi}) to confession of Jesus as Lord in
the Holy Spirit. The opposed formulas, ``Anathema to Jesus'' and ``The Lord
Jesus,'' provide a basic Christological criterion; they are not presented as
an exhaustive test that authenticates every subsequent spiritual claim.
Paul's own argument continues through the testing, intelligibility, mutual
edification, and order required in chapter~14.

Verses 4--6 coordinate three real differences without dividing their divine
source: distributions of gifts (\latin{diaireseis charismatōn}),
ministries (\latin{diakoniōn}), and operations
(\latin{energēmatōn}) stand under the same Spirit, the same Lord, and the
same God. The triadic form is neither a list of three unrelated economies nor
a warrant to collapse gift, ministry, and operation into synonyms. Verse~7
states the social end before the examples appear: manifestation is given
\latin{pros to sympheron}, toward benefit. Verse~11 then preserves divine
freedom and personal distribution: the one Spirit works all these things,
allotting individually to each as he wills
(\latin{idia hekastō kathōs bouletai}).

The canonical setting prevents the catalogue from becoming a prestige scale.
Paul has already treated idol food and participation in idolatrous worship
(1~Cor.~8:1--13; 10:14--22). After this list, one body has many mutually
dependent members (12:12--27), God appoints differentiated services and gifts
without giving all of them to all (12:28--31), love is the indispensable
``more excellent way'' (ch.~13), and intelligible speech, testing, edification,
and peace govern the assembly (ch.~14). Diversity is therefore neither
competition nor formless spontaneity: a manifestation is received, tested,
and exercised for the body.

John Chrysostom's \emph{Homily 29} supplies a running ancient reception.
Sections~1--3 contrast the former pagan compulsion he describes with rational
Christian confession; section~4 calls the manifestations gifts rather than
debts and reads verses 4--7 toward common profit; section~6 emphasizes that
the Spirit distributes and acts ``as he wills.'' Chrysostom's historical
account of baptismal manifestations and of gifts later becoming obscure is
his fourth-century reception, not a universal doctrinal settlement that all
such gifts ceased. Ambrosiaster, at 1~Corinthians 12:1--11 (PL~17,
cols.~150--152), likewise excludes boasting by grace and relates differentiated
gifts to one divine working and ecclesial usefulness.

Aquinas, commenting on chapter~12, preserves the distinctions among gifts,
ministries, and operations while returning them to one source and useful end.
Basil's \emph{On the Holy Spirit} 16.37 is a doctrinal reuse rather than a
running commentary: he cites verses 4--6 and 11 to confess inseparable divine
operation together with the Spirit's genuine agency. Vatican~II receives
verse~7 in \emph{Lumen gentium} 12 and \emph{Apostolicam actuositatem} 3:
charisms, whether striking or ordinary, are ordered to the Church's renewal
and building and remain subject to discernment by legitimate pastors. This
later vocabulary must not be retrojected as though it exhausted Paul's
first-century terms, but it rules out both private possession and untested
enthusiasm.

\begin{fourstable}{Textual movement}{Exegetical force}{Ecclesial consequence}{Controlling limit}
\textbf{Former leading} & Mute idols give way to confession of Jesus as Lord in the Spirit. & Discernment begins from Christological allegiance. & The formula is not by itself an exhaustive authentication test.\\
\textbf{Differences} & Gifts, ministries, and operations remain distinguishable under one divine source. & Difference can serve communion without becoming rivalry. & Charism, office, talent, and claimed phenomenon are not interchangeable terms.\\
\textbf{Manifestation} & Each manifestation is given toward benefit. & The recipient is accountable to the body rather than made its superior. & Spectacle, intensity, or self-report does not prove divine origin.\\
\textbf{Distribution} & The Spirit allots individually as he wills. & No member authors the gift's status, and no member possesses every gift. & Divine freedom does not remove testing, stewardship, pastoral judgment, or order.\\
\end{fourstable}
